\scanpage{331}

est muni canoniquement (non d'une quasi-rigidification discrète \inserted{orientée}, ce qui dépend du choix d'un isom $\overline{\mathbb Q}\simeq\overline{\mathbb Q}_0$ i.e. d'un plongement $\overline{\mathbb Q}\subset\mathbb C$) mais du moins d'un groupe de Teichmüller géométrique strict $\Sigma\subset\widehat{\widehat{\mathcal T}}(\pi)$, à savoir l'image de $\pi_1(M_{g,\nu\overline{\mathbb Q}})$, \inserted{isomorphe à} Ceci posé, On a un hom canonique d'extensions de groupes

\diagram{(26)\quad
\begin{tikzcd}[ampersand replacement=\&, column sep=small]
1 \arrow[r] \&
\pi_1(M_{g,\nu\overline{\mathbb Q}}) \arrow[r] \arrow[d, "{\unsure{\wr|2}}"] \&
\pi_1(M_{g,\nu}) \arrow[r, no head] \arrow[d] \&
\operatorname{Gal}(\overline{\mathbb Q}/\mathbb Q) \arrow[r] \arrow[d] \&
1 \\
1 \arrow[r] \&
\Sigma \arrow[r] \&
\unsure{\mathcal N|N}_\Sigma \arrow[r] \arrow[d, equal] \&
\Gamma_\Sigma \arrow[r] \arrow[dr] \&
1 \\
\& \&
\operatorname{Norm}_{\widehat{\widehat{\mathcal T}}(\pi)}(\Sigma) \&
\&
\text{\margin{groupe des autom. \deleted{\unsure{``arithmétiques}}\inserted{extérieurs} "arithmétiques" de $\pi$ (muni de $\Sigma$)}}
\end{tikzcd}}

Pour le choix d'un isom $\overline{\mathbb Q}\simeq\overline{\mathbb Q}_0$ i.e. d'un plongement $\overline{\mathbb Q}\subset\mathbb C$, l'image de la conjugaison complexe de $\operatorname{Gal}(\overline{\mathbb Q}_0/\mathbb Q)$ n'est autre que \unsure{l'élt} d'ordre $2$ de $\Gamma_\Sigma$ associé au sous-groupe de Teichmüller géométrique, associé à la quasi-discrétification (orientée, mais peu importe) canonique sur $\pi$ (définie par $\overline{\mathbb Q}\unsure{\simeq|\xrightarrow{\sim}}\overline{\mathbb Q}_0$).

\scanpage{332}
\mspage{168}

La conjecture naturelle ici, c'est que
\[
(27)\qquad \operatorname{Gal}(\overline{\mathbb Q}/\mathbb Q)\longrightarrow\Gamma_\Sigma
\]
du groupe de Galois vers les autom ext. arithmétiques de $(\pi,\Sigma)$ soit un iso. On peut la compléter par la conjecture, encore plus hardie, que de plus $\Sigma$ est invariant dans $\widehat{\widehat{\mathcal T}}(\pi)$ ($\Longleftrightarrow\Sigma_{g,\nu}$ inv. dans $\widehat{\widehat{\mathcal T}}_{g,\nu}$), ce qui signifierait donc que $\unsure{\mathcal N|N}_\Sigma=\widehat{\widehat{\mathcal T}}$, ou encore que \deleted{$\unsure{\pi_1(M_{g,\nu})}$}

\margin{Conjectural!!}
\[
(28)\qquad \pi_1(M_{g,\nu})\xrightarrow{\sim}\widehat{\widehat{\mathcal T}}(\pi)
=\operatorname{Aut\,ext}_{\mathrm{lac}}(\pi).
\]
et $\operatorname{Gal}(\overline{\mathbb Q}/\mathbb Q)$ s'identifierait au quotient "arithmétique" du groupe des autom. ext. à lacets de $\pi$.

Soit maintenant $U$ une courbe algébrique de type $g,\nu$ sur $\overline{\mathbb Q}$. Elle est définie par un pt de $M_{g,\nu\overline{\mathbb Q}}$ sur $\operatorname{Spec}\overline{\mathbb Q}$, et \deleted{choisissant un} comme fibre de $\mathcal U_{g,\nu}$ en ce pt. Choisissons un rev. universel de $\unsure{U}$, d'où un groupe $\pi(U)$, canoniquement isomorphe au groupe $\pi$ ci-

\scanpage{333}

dessus, une fois choisi un isom entre $\widetilde{\mathcal U}_{g,\nu}$ et le rev. universel de $\mathcal U_{g,\nu}$ déduit de $\widetilde U$, ce qui donne une indétermination dans $\pi_1(\mathcal U_{g,\nu\overline{\mathbb Q}})\simeq\unsure{\widehat{\mathcal S}^{+}}$.

Ainsi, une fois choisi un isom $\overline{\mathbb Q}\simeq\overline{\mathbb Q}_0$, on trouve sur $\pi_1(U)$ une quasi-discrétification orientée (évidente d'ailleurs à définir \unsure{directt}, par voie \unsure{transcend.te}) changée en son opposée par chgt de l'isom $\unsure{\overline{\mathbb Q}_2|\overline{\mathbb Q}}\simeq\overline{\mathbb Q}_0$ par conj. complexe. D'autre part, le sous groupe $\widehat{\mathcal T}^{+}$ de $\widehat{\widehat{\mathcal T}}(\pi)$ correspondant à cette quasi-discrétification (orientée, peu nous chaut) ne dépend pas du choix de l'isom en question $\overline{\mathbb Q}\simeq\overline{\mathbb Q}_0$, il est canoniquement associé à $\pi$ en tant que groupe fond. d'une $U$ sur \deleted{$\unsure{\overline{\mathbb Q}}$} un corps alg clos

[\uline{NB} on pourrait prendre un corps alg clos de car. $0$ quelconque, pas la peine que ce soit \inserted{\unsure{un|sur}} $\overline{\mathbb Q}$, cf plus bas].

\scanpage{334}
\mspage{169}

\deleted{Considérons maintenant une courbe algébrique $U$ de type $(g,\nu)$ définie sur un corps $K$ de type fini --- donc définie par un morphisme de $\mathbb Q$-multiplicités
\[
\operatorname{Spec}K\mathrel{\text{---}}M_{g,\nu}
\]
comme image inverse de $\mathcal U_{g,\nu}$. On a donc canoniquement un \unsure{hom} cartésien groupoïdes fondamentaux
\diagram{\begin{tikzcd}[ampersand replacement=\&]
\Pi_1(U) \arrow[r] \arrow[d] \& \Pi_1(\mathcal U_{g,\nu}) \arrow[d] \\
\Pi_1(K) \arrow[r,dash] \& \Pi_1(M_{g,\nu})
\end{tikzcd}}
}

Si la conjecture sur $\Gamma$ est vérifiée, il en résulterait que la donnée d'un groupe \deleted{\unsure{extérieur}} profini à lacets \inserted{$\pi$} de type $(g,\nu)$, muni d'un groupe de Teichmüller géom. strict $\Sigma\subset$ $\widehat{\widehat{\mathcal T}}(\pi)$ (ce qui peut être n'est pas une structure supplémentaire du tout --- $\Sigma$ serait uniquement déterminé par $\pi$ !) : équivaudrait à la donnée $1^{\circ})$ d'une extension \unsure{alg} close $\overline{\mathbb Q}$ de $\mathbb Q$ (dont le groupe de Galois serait $N_{\widehat{\widehat{\mathcal T}}}(\Sigma)/\Sigma=\Gamma_\Sigma$, et le \inserted{\unsure{$\Gamma_0$-|$\Pi_0$-}}torseur $\operatorname{Isom}(\overline{\mathbb Q}_0,\overline{\mathbb Q}_1)$ s'identifierait

\scanpage{335}

à l'ens \inserted{($\subset\operatorname{Isom}(\widehat{\Pi}_{g,\nu},\pi)/\widehat{\mathcal T}_{g,\nu}^{\,0}$} des quasi-rigidifications \deleted{\illegible{1}} \uline{orientées} donnant naissance à \unsure{$\pi_{g,\nu}$|$\Pi_{g,\nu}$}) et d'un rev. universel de $\mathcal U_{g,\nu\,\overline{\mathbb Q}}$ --- ou \unsure{simplt}, d'un rev. universel $\widetilde{\mathcal U}_{g,\nu}$ de $\mathcal U_{g,\nu}=\mathcal U_{g,\nu\,\mathbb Q}$.

\margin{Si on se donne \unsure{seulement} $\pi$ comme groupe ext., avec $\Sigma$, ceci revient à la donnée d'un rev. universel de $M_{g,\nu}$.}

La donnée d'un isom $\overline{\mathbb Q}\simeq\overline{\mathbb Q}_0$ revient donc à celle d'une quasi-rigidification orientée sur $\pi$ compatible \unsure{avec} $\Sigma$ (condition peut être automatiquement satisfaite\ldots).
\margin{\uline{NB} Si on renonce à choisir $\widetilde U$, ceci \unsure{équivaut|revient} à travailler \unsure{avec} des groupes à lacets \uline{extérieurs}.}
Ceci dit, la donnée d'une courbe algébrique \inserted{de type $g,\nu$} sur \inserted{\deleted{et d'un rev. universel}} \deleted{\unsure{(}\illegible{1} \unsure{correspondrait}} un corps alg. clos $\overline{\mathbb Q}$, et d'un rev. universel \deleted{de celle-ci} de celle ci reviendrait à la donnée d'une donnée précédente (\unsure{savoir} \inserted{\deleted{un \unsure{groupe ext. profini}} un rev. universel d'un $M_{g,\nu}$}), \uline{plus} un pt de $M_{g,\nu,\overline{\mathbb Q}}$ sur $\overline{\mathbb Q}$ \deleted{\unsure{cette}} \inserted{\unsure{définissant ce rev. univ.} --- et,} cette dernière donnée s'identifierait à un \deleted{\unsure{point}} germe de scindage de l'extension
\inserted{\[
1\longrightarrow\pi_1(M_{g,\nu\,\overline{\mathbb Q}})\longrightarrow\pi_1(M_{g,\nu})\longrightarrow\pi_1(\mathbb Q)\longrightarrow1
\]}
\unsure{L'}interprétation \inserted{profinie} \inserted{(en termes des \unsure{groupes} \uline{ext.} à lacets $\pi$)} \inserted{\unsure{est alors}} \deleted{\unsure{$\longrightarrow N_\Sigma\longrightarrow\Gamma_\Sigma\longrightarrow1$}} un $(\pi,\Sigma)$, et un germe de scindage de
\[
1\longrightarrow\Sigma\longrightarrow N_\Sigma\longrightarrow\Gamma_\Sigma\longrightarrow1.
\]
Si on veut une courbe sur une sous-extension finie \unsure{$K'$} de $\overline{\mathbb Q}/\mathbb Q$, correspondant à un ss-groupe d'indice fini $\Gamma'\subset\Gamma_\Sigma$, il s'agira d'un scindage \deleted{\illegible{1}} partiel $\Gamma'\longrightarrow N_\Sigma$.

\scanpage{336}
\mspage{170}

Mais sûrement, si cette \deleted{\unsure{représentation}} \inserted{description} des courbes algébriques anabéliennes est pleinement fidèle, elle n'est pas $2$-fidèle, i.e. il faut des conditions sur un scindage, qu'il faudra examiner par la suite.

La description d'une \deleted{\unsure{gr.}} courbe alg. de type $g,\nu$ sur un corps \inserted{fixé} de type fini \inserted{\unsure{, car.\ $0$|sur $\mathbb Q$}} $K$ quelconque, muni d'une extension alg close $\overline K$ (donc \unsure{$\Gamma=\operatorname{Gal}\overline K/K'$} $\longrightarrow\Gamma_0=\operatorname{Gal}(\overline{\mathbb Q}/\mathbb Q)$, où $\overline{\mathbb Q}$ est la clôture algébrique de $\mathbb Q$ dans $\overline K$), en termes de $\Gamma\longrightarrow\Gamma_0$, serait la suivante : donnée d'un groupe extérieur \unsure{à} lacets $\pi$ de type $g,\nu$, d'un \unsure{ss-groupe} de Teichmüller $\Sigma$ dans $\widehat{\widehat{\mathcal T}}(\pi)$, d'un isom $\overline{\mathbb Q}\simeq$ corps défini par cette situation (ayant comme groupe de Galois \deleted{par} $N_\Sigma/\Sigma$, de sorte qu'on trouve un iso. $\Gamma_0\simeq N_\Sigma/\Sigma$) enfin d'un relèvement de $\Gamma\longrightarrow\Gamma_0$ \deleted{\illegible{2}} \inserted{(qui décrit} une action arithmétiquement extérieure de $\Gamma$ sur $(\pi,\Sigma)$) en une action extérieure $\Gamma\longrightarrow N_\Sigma$ de $\Gamma$ sur $(\pi,\Sigma)$ ; avec des conditions qu'il faudra \unsure{essayer}

\scanpage{337}
de dégager (néc. et suff. conjecturalement) sur ce relèvement.

\deleted{Pour que}
Notons que tout plongement $\overline K\subset\mathbb C$
définit \inserted{\unsure{canont|canoniq.}} sur le groupe extérieur
$\pi=\pi_1(U_{\overline K})$
une quasi-discrétification orientée, \unsure{remplacée}
par l'opposée quand on remplace le plongement par le complexe conjugué.
\deleted{Supposons} \deleted{\unsure{qu'|que}} Je dis que ces
quasi-discrétifications définissent toutes $\Sigma$ comme groupe des
\unsure{autom.}, et que la quasi-discrétification
définie par $\overline K\hookrightarrow\mathbb C$ ne dépend que de
sa restriction en
$\overline{\mathbb Q}\subset\mathbb C$
--- de \unsure{façon} plus précise, c'est celle qu'on définit
\unsure{directt} à l'aide de $\overline{\mathbb Q}\subset\mathbb C$ i.e. de
$\overline{\mathbb Q}\simeq\overline{\mathbb Q}_0$
--- Mais il y a lieu d'examiner aussi la façon d'obtenir des
\uline{discrétifications}. Ainsi, si \inserted{\unsure{on a un|un}} $U$ est définie
sur $\overline{\mathbb Q}_0$ (i.e. un \unsure{$\pi$|$\Pi$} ext. de type
$(g,\nu)$, muni d'une quasi-discrétification

\scanpage{338}
\mspage{171}
orientée \unsure{$\pi_0^{\natural+}$|$\pi_0^{\zeta+}$},
et d'un germe de \deleted{\unsure{relèvement}}\inserted{scindage}
de
\diagram{
\begin{tikzcd}[ampersand replacement=\&, column sep=large]
1 \arrow[r] \&
\Sigma_{\deleted{\unsure{\pi_0^{\natural+}}}}
\arrow[r] \arrow[d, equal] \&
N_{\Sigma} \arrow[r] \& \Gamma_{\Sigma} \arrow[r] \& 1 \\
\& \mathrm{Aut\,ext}_{\unsure{\mathrm{lac}}}
 (\pi,\unsure{\pi_0^{\natural+}|\pi_0^{\zeta+}}) \& \& \&
\end{tikzcd}
}
\unsure{)|,}
qui fait opérer \inserted{ext.} le noyau \unsure{de|du} groupe
défini par $\Gamma_{\Sigma}$ sur $\pi$, les pts rat/$\overline{\mathbb Q}_0$
\unsure{correspondants|correspondent} aux
\inserted{classes de $\pi$-conj. des} germes de relèvement de
cette opération\unsure{,} \unsure{ext.|ext} en une vraie opération.
Supposons donc donné un tel point, i.e.
\deleted{\unsure{$\Gamma_{\Sigma}$ opère}} \unsure{on a} un sous-groupe ouvert
\unsure{$\Gamma^{\pi+}$|$\Gamma'$}$\subset\Gamma_{\Sigma}$
qui opère bel et bien sur $\pi$,
\inserted{l'op. définie} mod autom intérieurs
(lui-$\overline{\mathrm m}$ \unsure{uniques}
car \unsure{$\pi^{\Gamma}$|$\overline\pi^{\Gamma}$}$=(1)$ !)
Dans cette situation, il faudrait définir dans $\pi$
\inserted{$(\simeq\pi_1(U_{\unsure{P|p}},P))$}
une \uline{discrétification}
\inserted{\uline{orientée} $\pi_0\subset\pi$,}
et pas seulement une quasi-discrétification orientée !
\deleted{Mais} (Bien sûr, elle doit être dans la classe \unsure{qu'on} \inserted{s'est} donnée d'avance
de \deleted{quasi} discrétifications orientées, \unsure{qu'on avait}
justement notée \unsure{$\pi_0^{\natural+}$|$\pi_0^{\zeta+}$}.
On y reviendra --- et à la situation analogue sur un corps de
base $K$ de type fini quelconque\unsure{....|...}

\scanpage{339}

Je \unsure{réfère} au début du \S\ suivant (\S\,27) pour le
changement de terminologie, la ``quasi-discrétification''
devenant une ``pré-discrétification''. Mais il y a lieu
d'introduire aussi une notion plus fine, correspondant
--- \unsure{au} \illegible{1} $\pi_1$ (profini) d'une \unsure{variété} alg.
$X$ définie sur un corps alg. clos $\overline K$, quand on plonge
$\overline K$ dans $\mathbb C$ : il y a une discrétification \unsure{mod}
\unsure{autom.|automs.} intérieurs --- on l'appellera une
prédiscrétification stricte. Dans le cas d'un
$\pi$ à lacets, l'\deleted{espace} homogène \deleted{de} dans $\widehat{\widehat{\mathcal T}}(\pi)$ de
celle ci s'identifie à
\[
\operatorname{Isom}_{\mathrm{lac}}(\widehat{\Pi}_{g,\nu},\pi)/\widehat{\Pi}_{g,\nu}\cdot
\overset{\deleted{\unsure{\widehat{\wedge}}}}{\mathcal T_{g,\nu}}
\inserted{\simeq\operatorname{Isom\,ext}_{\mathrm{lac}}(\widehat{\Pi}_{g,\nu},\pi)/\mathcal T_{g,\nu}}
\]
--- cette notion \inserted{\unsure{égalt|égalmt}} ne dépend que du gr. \unsure{pr} à lacets
extérieur défini par $\Pi_{g,\nu}$), et \deleted{\illegible{1}} \unsure{équivaut}
à celle d'une ``base extérieure'' dans $\pi$
\inserted{\unsure{mod action de} $\mathcal T_{g,\nu}$,}
\inserted{i.e. d'un iso.}
de $\pi$ avec le $\widehat{\Pi}_{g,\nu}$ ``type'' ;
\inserted{\unsure{mod. action de} $\mathcal T_{g,\nu}$.}
L'application
\[
\overset{\inserted{\text{Bases extérieures de }\pi}}{\deleted{\text{Prédiscrétifications strictes }(\pi)}}
\longrightarrow \text{Prédiscrétifications de }(\pi)
\]
fait \unsure{donc} du membre de gauche un torseur relatif
\inserted{à droite,} sur celui de droite, du groupe $\mathcal T_{g,\nu}$.
\deleted{\illegible{1}}

Une prédiscrétification de $(\pi,\Sigma)$ correspond à un
choix d'un isom de la clôture algébrique $\overline{\mathbb Q}_{\unsure{(\pi,\Sigma)}}$ \inserted{$=\overline{\mathbb Q}$} \unsure{associée}
à $(\pi,\Sigma)$, avec $\overline{\mathbb Q}_0$. Quand on a donné $(\pi,\Sigma)$ comme
\unsure{correspondant} à une courbe algébrique sur $\overline{\mathbb Q}$, i.e. qu'on
\unsure{se donne|s'est donné} un \unsure{germe} de relèvement de $\Gamma_\Sigma$ dans $N_\Sigma\subset\widehat{\widehat{\mathcal T}}(\pi)$
i.e. un \unsure{germe|groupe} d'action extérieure de $\Gamma_\Sigma$ sur $\pi$, \unsure{toute|dont}
prédiscrétification \inserted{doit} donner naissance à une discrétification
stricte --- et \inserted{une} \deleted{\illegible{1}} discrétification \unsure{équivaut} \unsure{en remontant}
un \unsure{germe} d'action (``admissible'') de $\Gamma_\Sigma$ sur $\pi$ \unsure{\ldots}

\scanpage{340}
\mspage{172}

27. \uline{Changement de type} $(g,\nu)$ : a) \uline{bouchage de trous}.

Je vais changer de terminologie, en appelant
pré-discrétification ce que j'avais appelé (un peu
péjorativement !) \unsure{quasi}-discrétification. En effet, on
\unsure{prévoit} qu'une pré-discrétification \inserted{--- compatible avec $\Sigma$ \unsure{---)}} (qui revient \unsure{moralement|moralt}
au plongement \inserted{(mod conj. complexe)} dans $\mathbb C$ de la clôture alg. $\overline{\mathbb Q}$ de $\mathbb Q$
canoniquement associée \unsure{à|$*$} $(\pi,\Sigma(\subset\widehat{\widehat{\mathcal T}}(\pi)))$, donne
naissance dès que l'action extérieure-arithmétique de
$\Gamma_\Sigma=\unsure{N_\Sigma/\Sigma}$ sur $\pi$ est relevée en un germe d'action
sur $\pi$, à une vraie discrétification de $\pi$. \unsure{Itou|Idem} pour
pré-discrétifications orientées. Le groupe $N_\Sigma/\Sigma$ mérite
\margin{\unsure{moralement|normalement}, $\Sigma=\unsure{\widehat{\mathcal T}}$ (mais il n'y a pas de $\unsure{\mathcal T}$ !)}
aussi un nom --- je vais l'appeler le \uline{groupe de}
\uline{Teichmüller} \uuline{arithmétique} associé à $(\pi,\Sigma)$, et
ses éléments seront appelés les automorphismes
arithmétiquement extérieurs de $\pi$ --- \unsure{déduits} d'un
autom extérieur (normalisant $\Sigma$) : en négligeant les
autom extérieurs ``géométriques'' (i.e. justement ceux
$\in\Sigma$) --- \deleted{\unsure{de}} comme les autom extérieurs \inserted{ordinaires} \unsure{étaient}
\unsure{décrits} en négligeant les autom intérieurs de $\pi$.
On \deleted{\illegible{1}} fera attention que le caractère ``multiplicateur''
\[
\chi:\widehat{\widehat{\mathcal T}}(\pi)\unsure{\longrightarrow}\widehat{\mathbb Z}^{*}
\]
est trivial sur $\Sigma$, par construction de $\Sigma$, et
passe à $\Gamma_\Sigma$
\[
\chi_\Sigma:\Gamma_\Sigma\longrightarrow\widehat{\mathbb Z}^{*}
\]
